% Carta de apresentação (cover letter) — submissão ao Bolema
% Compilação: tectonic calculo-bolema-carta.tex
\documentclass[12pt,a4paper,oneside]{article}
\usepackage[T1]{fontenc}
\usepackage[utf8]{inputenc}
\usepackage[brazilian]{babel}
\usepackage{newtxtext}
\usepackage[a4paper,top=2cm,bottom=2cm,left=2.8cm,right=2.3cm]{geometry}
\usepackage{setspace}
\usepackage[hidelinks]{hyperref}
\usepackage{parskip}
\usepackage{enumitem}
\setlist[itemize]{leftmargin=1.3em,itemsep=1pt,topsep=1pt}
\setstretch{1.06}
\setlength{\parskip}{4pt}
\pagestyle{empty}

\begin{document}

\noindent\hfill 2026.

\bigskip
\noindent Aos Editores-chefes do \textbf{Bolema: Boletim de Educação Matemática}\\
Prof.\ Dr.\ Danyal Farsani \textbullet{} Prof.\ Dr.\ Marcus Vinicius Maltempi \textbullet{} Prof.\ Dr.\ Roger Miarka\\
Universidade Estadual Paulista (Unesp), Rio Claro --- SP

\bigskip
\noindent Prezados Editores,

Temos a satisfação de submeter à apreciação do \textit{Bolema}, para
avaliação na seção de \textbf{Artigos}, o manuscrito intitulado \textbf{``O
cálculo ausente: duzentos anos de currículo, uma comparação internacional e
um roteiro pedagógico para o ensino médio''}, de nossa autoria e caráter
original e inédito.

\medskip
\noindent\textbf{Do que trata o artigo.}
O trabalho conduz uma revisão documental comparada dos currículos oficiais de
matemática do ensino médio, com foco específico na presença ou ausência de
cálculo diferencial e integral antes do ingresso no ensino superior. A partir
dos documentos curriculares oficiais de dez países de cinco continentes --- de
alta, média e baixa renda --- e do programa \textit{International
Baccalaureate}, o artigo documenta um achado que consideramos relevante para a
comunidade de Educação Matemática: \textbf{o Brasil é o único país da amostra
cujo currículo nacional do ensino médio exclui, simultânea e integralmente, o
cálculo diferencial e o integral}. Vários países com Índice de Desenvolvimento
Humano inferior ao brasileiro (Índia, Vietnã, Paquistão, Nigéria) oferecem
cálculo completo --- o que afasta a explicação por renda e devolve a questão ao
terreno curricular e pedagógico.

\medskip
\noindent\textbf{Por que o \textit{Bolema}.}
Entendemos que o manuscrito se inscreve no escopo central da revista --- o
ensino e a aprendizagem de Matemática e o papel da Educação Matemática na
sociedade --- e dialoga com a tradição de estudos curriculares que o próprio
\textit{Bolema} consolidou: da edição temática dedicada ao currículo de
matemática (Pires; Godoy; Silva; Santos, 2014) aos estudos comparados entre
sistemas nacionais (Gonçalves; Dias; Peralta, 2018) e à discussão sobre a
introdução das ideias do cálculo na educação básica (Araújo; Avelar, 2022).
Diferentemente desses trabalhos, nosso foco é estritamente a presença ou
ausência do cálculo no currículo prescrito, mapeada país a país.

O artigo não se limita, porém, ao diagnóstico. Ancorado no currículo em
espiral de Bruner, na zona de desenvolvimento proximal de Vygotsky e no
diagnóstico epistemológico de Rezende --- e articulado à literatura
internacional da didática do cálculo (imagem e definição conceitual em Tall e
Vinner; reificação em Sfard; a teoria APOS; as dificuldades da análise em
Artigue) ---, ele lê a ausência brasileira à luz da \textbf{transposição
didática} de Chevallard: uma \textit{não-transposição} do saber de referência
ao saber a ensinar. E oferece, na tradição da \textbf{engenharia didática}, o
desenho e a análise \textit{a priori} de uma sequência para introduzir derivada
e integral no ensino médio segundo a lógica \textit{Concrete-Pictorial-Abstract}
--- explicitando que sua validação empírica em sala constitui a etapa seguinte
da pesquisa.

\medskip
\noindent\textbf{Contribuições que gostaríamos de destacar.}
\begin{itemize}
  \item \textbf{Um mapeamento inédito e preciso} de um contraste internacional
  pouco quantificado na literatura nacional, com formulação cuidadosa dos
  limites de inferência: o desempenho brasileiro em avaliações internacionais e
  as taxas de reprovação em Cálculo~I são apresentados como indicadores
  contextuais, não como prova causal, e a correlação entre status curricular e
  PISA é reportada com honestidade, inclusive quando nula.
  \item \textbf{Reprodutibilidade auditável}, rara neste gênero: o achado
  central é verificável por terceiros em minutos, mediante um manifesto público
  das fontes curriculares oficiais (com \textit{hash} SHA-256 e data de captura)
  e roteiros que extraem e conferem os documentos. Para a avaliação cega, os
  artefatos estão disponíveis em espelho anônimo.
  \item \textbf{Uma ponte para a sala de aula}: a sequência em engenharia
  didática converte o diagnóstico em proposta fundamentada, à disposição da
  comunidade.
\end{itemize}

\medskip
\noindent\textbf{Declarações.}
\begin{itemize}
  \item O artigo é original e inédito, não publicado nem submetido
  simultaneamente a outro periódico.
  \item Não há conflito de interesse e a pesquisa não recebeu financiamento externo.
  \item Todos os autores aprovaram a versão submetida e por ela se responsabilizam.
  \item Declaramos o uso do assistente de inteligência artificial Claude (Opus
  4.8, Anthropic) como ferramenta de apoio em todas as etapas (curadoria e
  revisão de fontes, redação e edição, geração do código das figuras), com
  revisão e verificação humanas integrais --- em atenção à política de
  transparência da revista, declaração também registrada no manuscrito.
  \item O arquivo do manuscrito foi anonimizado (sem nomes, afiliações ou
  identificadores no corpo do texto nem nas propriedades do arquivo), para a
  avaliação por pares duplo-cega.
\end{itemize}

\medskip
Colocamo-nos à disposição para quaisquer esclarecimentos e agradecemos a
atenção dispensada.

\bigskip
\noindent Atenciosamente,

\medskip
\noindent\textbf{Leonardo Chalhoub} \, (autor correspondente)\\
Mestre em Administração (Finanças), UFRGS --- pesquisador independente, Brasil\\
E-mail: \href{mailto:leochalhoub@hotmail.com}{leochalhoub@hotmail.com} \textbullet{}
ORCID: 0000-0003-0484-158X \textbullet{}
Lattes: \url{http://lattes.cnpq.br/0703595867512327}

\medskip
\noindent\textbf{Coautores}
\begin{itemize}
  \item \textbf{Alexandre Maciel Rolim} --- Mestre em Proteção Radiológica
  (IFSC), Florianópolis--SC \textbullet{}
  \href{mailto:alexandremrolim@gmail.com}{alexandremrolim@gmail.com} \textbullet{}
  ORCID: 0009-0000-3983-4713 \textbullet{} Lattes: \url{http://lattes.cnpq.br/0128707599536129}
  \item \textbf{Luis Fernando Kranz} --- Mestre em Administração (UFRGS),
  Porto Alegre--RS \textbullet{}
  \href{mailto:lufekranz@gmail.com}{lufekranz@gmail.com} \textbullet{}
  ORCID: 0009-0009-6015-3141 \textbullet{} Lattes: \url{http://lattes.cnpq.br/9302324694749099}
  \item \textbf{Jefferson Korte Junior} --- graduando em Ciência da Computação
  (UTFPR), Curitiba--PR \textbullet{}
  \href{mailto:jefferson.2024@alunos.utfpr.br}{jefferson.2024@alunos.utfpr.br} \textbullet{}
  ORCID: 0009-0006-9466-9830 \textbullet{} Lattes: \url{http://lattes.cnpq.br/6748323429576556}
\end{itemize}

\end{document}
